\subsection{Correlazione}\label{sec:methods_correlation}

%Ci sono fattori sistematici che influenzano la performance dei vari Paesi alle
%Olimpiadi?

% info generali e probabilmente scontate, le mettiamo in generale solo se
% entrano

%Nella letteratura esistente, due dei fattori che più vengono tenuti in
%considerazione sono la ricchezza e la popolazione di un Paese: \begin{itemize}
%\item la ricchezza, in quanto permette al Paese di investire nelle
%infrastrutture sportive e nel sostegno agli atleti \item la popolazione, in
%quanto al crescere del numero di abitanti, probabilisticamente, cresce la
%dimensione del bacino di potenziali talenti \end{itemize}
%
%Per verificare correlazione tra ricchezza/popolazione e performance Olimpica,
%abbiamo deciso di utilizzare la correlazione di Spearman.

% altre descrizioni implementative e spiegazioni non necessarie qui

%Il coefficiente di correlazione di Spearman è una misura statistica non
%parametrica che valuta la forza e la direzione della relazione tra due
%variabili. A differenza della correlazione di Pearson, che misura la relazione
%lineare, Spearman valuta la relazione monotonica.
%
%Il calcolo non si basa sui valori grezzi dei dati, ma sui loro ranghi (ovvero
%sulla loro posizione in classifica). Il procedimento prevede di:
%
%\begin{enumerate} \item Ordinare i dati di ciascuna variabile dal più piccolo
%    al più grande. \item Assegnare un rango (1, 2, 3...) a ogni valore. \item
%    Calcolare il coefficiente di correlazione basandosi sulla differenza tra i
%    ranghi delle due variabili. \end{enumerate}

% fonti per queste cose catturate? soprattutto se in dubbio, se queste cose non
% sono state determinanti nei risultati, si possono togliere.

%Abbiamo deciso di utilizzare Spearman invece del più classico Pearson per due
%ragioni principali: \begin{itemize} \item Spearman non richiede che i dati
%seguano una distribuzione normale \item Spearman è in grado di catturare
%relazioni che Pearson ignorerebbe, come ad esempio una crescita esponenziale o
%logaritmica. Spearman non richiede che i dati crescano secondo una linea retta;
%è sufficiente che, al crescere di una variabile, l'altra tenda a crescere (o
%decrescere) in modo costante, anche se con un ritmo non uniforme. \end{itemize}

% unito nei risultati

%Il coefficiente di Spearman varia da -1 (perfetta relazione monotonica
%negativa) a +1 (perfetta relazione monotonica positiva). Uno 0 indica assenza
%di relazione.


Per mostrare correlazione tra ricchezza/popolazione e performance Olimpica,
abbiamo deciso di utilizzare la correlazione di Spearman, che valuta la
relazione monotonica tra due variabili.\\
Pearson avrebbe richiesto una distribuzione normale dei dati.

% ma questo requisito, a noi vale?

%    \item Spearman è in grado di catturare relazioni che Pearson ignorerebbe,
%    come ad esempio una crescita esponenziale o logaritmica. Spearman non
%    richiede che i dati crescano secondo una linea retta; è sufficiente che, al
%    crescere di una variabile, l'altra tenda a crescere (o decrescere) in modo
%    costante, anche se con un ritmo non uniforme.
